\section*{TDK Abstract Submission}
\addcontentsline{toc}{section}{TDK Abstract}

\textbf{Title:} PhishGuard: OSINT-Enhanced Machine Learning Architecture for Real-Time Phishing Detection

\textbf{Author:} Ishaq Muhammad (PXPRGK)

\textbf{Supervisor:} Md. Easin Arafat

\textbf{Section:} Information Technology / Computer Science

\vspace{1em}

\textbf{Abstract (maximum 10,000 characters):}

Phishing attacks remain the predominant cybersecurity threat, causing over \$10 billion in annual losses globally. Traditional reactive defense mechanisms---such as URL blacklists---fail against modern ephemeral campaigns where 50\% of phishing domains exist for less than 12 hours. This research presents \textit{PhishGuard}, a hybrid threat intelligence platform that orchestrates machine learning (ML), natural language processing (NLP), and Open-Source Intelligence (OSINT) to detect zero-day phishing attacks in real-time.

\textbf{Key Contributions:}

1. \textbf{OSINT-Enhanced ML Pipeline:} Unlike purely lexical classifiers, PhishGuard integrates 21 structural features with 4 real-time OSINT signals (WHOIS domain age, DNS infrastructure, VirusTotal reputation, AbuseIPDB reports). This contextual awareness enables detection of sophisticated threats that bypass structural analysis alone.

2. \textbf{Asynchronous Orchestration:} The system implements a FastAPI backend with async/await patterns for concurrent OSINT queries, achieving sub-3-second analysis latency while maintaining API responsiveness under load.

3. \textbf{Explainable AI (XAI):} Beyond binary classification, PhishGuard synthesizes empirical scores into human-readable explanations using SHAP values and heuristic mapping, educating users about detected threat indicators.

4. \textbf{Production Deployment:} The complete system is deployed as a publicly accessible web application (Next.js frontend on Vercel, Python backend on Render), demonstrating real-world viability.

\textbf{Empirical Results:}

An XGBoost classifier trained on 23,374 labeled samples achieved 96.45\% accuracy, 97.86\% precision, and 99.41\% AUC-ROC on a held-out test set (5,009 samples). Ablation studies demonstrate that OSINT features contribute significantly to detection of zero-day threats missed by lexical-only approaches.

\textbf{Keywords:} phishing detection, OSINT, machine learning, XGBoost, cybersecurity, natural language processing, explainable AI

\textbf{Repository:} \url{https://github.com/ishaq2321/phishing-detection-osint}

\textbf{Live Demo:} \url{https://project-4soy4.vercel.app}
